% Options for packages loaded elsewhere
\PassOptionsToPackage{unicode}{hyperref}
\PassOptionsToPackage{hyphens}{url}
%
\documentclass[
]{article}
\usepackage{amsmath,amssymb}
\usepackage{iftex}
\ifPDFTeX
  \usepackage[T1]{fontenc}
  \usepackage[utf8]{inputenc}
  \usepackage{textcomp} % provide euro and other symbols
\else % if luatex or xetex
  \usepackage{unicode-math} % this also loads fontspec
  \defaultfontfeatures{Scale=MatchLowercase}
  \defaultfontfeatures[\rmfamily]{Ligatures=TeX,Scale=1}
\fi
\usepackage{lmodern}
\ifPDFTeX\else
  % xetex/luatex font selection
\fi
% Use upquote if available, for straight quotes in verbatim environments
\IfFileExists{upquote.sty}{\usepackage{upquote}}{}
\IfFileExists{microtype.sty}{% use microtype if available
  \usepackage[]{microtype}
  \UseMicrotypeSet[protrusion]{basicmath} % disable protrusion for tt fonts
}{}
\makeatletter
\@ifundefined{KOMAClassName}{% if non-KOMA class
  \IfFileExists{parskip.sty}{%
    \usepackage{parskip}
  }{% else
    \setlength{\parindent}{0pt}
    \setlength{\parskip}{6pt plus 2pt minus 1pt}}
}{% if KOMA class
  \KOMAoptions{parskip=half}}
\makeatother
\usepackage{xcolor}
\usepackage[margin=1in]{geometry}
\usepackage{longtable,booktabs,array}
\usepackage{calc} % for calculating minipage widths
% Correct order of tables after \paragraph or \subparagraph
\usepackage{etoolbox}
\makeatletter
\patchcmd\longtable{\par}{\if@noskipsec\mbox{}\fi\par}{}{}
\makeatother
% Allow footnotes in longtable head/foot
\IfFileExists{footnotehyper.sty}{\usepackage{footnotehyper}}{\usepackage{footnote}}
\makesavenoteenv{longtable}
\usepackage{graphicx}
\makeatletter
\def\maxwidth{\ifdim\Gin@nat@width>\linewidth\linewidth\else\Gin@nat@width\fi}
\def\maxheight{\ifdim\Gin@nat@height>\textheight\textheight\else\Gin@nat@height\fi}
\makeatother
% Scale images if necessary, so that they will not overflow the page
% margins by default, and it is still possible to overwrite the defaults
% using explicit options in \includegraphics[width, height, ...]{}
\setkeys{Gin}{width=\maxwidth,height=\maxheight,keepaspectratio}
% Set default figure placement to htbp
\makeatletter
\def\fps@figure{htbp}
\makeatother
\setlength{\emergencystretch}{3em} % prevent overfull lines
\providecommand{\tightlist}{%
  \setlength{\itemsep}{0pt}\setlength{\parskip}{0pt}}
\setcounter{secnumdepth}{-\maxdimen} % remove section numbering
\usepackage[spanish]{babel}
\usepackage{amsmath}
\usepackage{fontspec}
\usepackage{unicode-math}
\usepackage{graphicx}
\usepackage{adjustbox}
\usepackage{tikz}
\usepackage{pgfplots}
\usepackage{booktabs}
\usetikzlibrary{3d,babel}
\ifLuaTeX
  \usepackage{selnolig}  % disable illegal ligatures
\fi
\usepackage{bookmark}
\IfFileExists{xurl.sty}{\usepackage{xurl}}{} % add URL line breaks if available
\urlstyle{same}
\hypersetup{
  hidelinks,
  pdfcreator={LaTeX via pandoc}}

\author{}
\date{\vspace{-2.5em}}

\begin{document}

\section{Question}\label{question}

En la gráfica se muestra la probabilidad de que la variable aleatoria
\(x\) tome valores en cada uno de tres intervalos en que se ha dividido
la curva.

\includegraphics[width=12cm,height=\textheight]{prob_dist_grafico.png}

¿Cuál de las siguientes tablas representa la probabilidad de que la
variable aleatoria \(x\) tome los valores en el intervalo indicado?

\subsection{Answerlist}\label{answerlist}

\begin{itemize}
\item
  \begin{longtable}[]{@{}cc@{}}
  \toprule\noalign{}
  Intervalo & Probabilidad \\
  \midrule\noalign{}
  \endhead
  \bottomrule\noalign{}
  \endlastfoot
  \(0 \le x \le 5\) & 0,41 \\
  \(5 < x \le 7\) & 0,30 \\
  \(7 < x \le 14\) & 0,30 \\
  \end{longtable}
\item
  \begin{longtable}[]{@{}cc@{}}
  \toprule\noalign{}
  Intervalo & Probabilidad \\
  \midrule\noalign{}
  \endhead
  \bottomrule\noalign{}
  \endlastfoot
  \(0 \le x \le 5\) & 0,30 \\
  \(5 < x \le 7\) & 0,71 \\
  \(7 < x \le 14\) & 1,00 \\
  \end{longtable}
\item
  \begin{longtable}[]{@{}cc@{}}
  \toprule\noalign{}
  Intervalo & Probabilidad \\
  \midrule\noalign{}
  \endhead
  \bottomrule\noalign{}
  \endlastfoot
  \(0 \le x \le 5\) & 0,30 \\
  \(5 < x \le 7\) & 0,41 \\
  \(7 < x \le 14\) & 0,30 \\
  \end{longtable}
\item
  \begin{longtable}[]{@{}cc@{}}
  \toprule\noalign{}
  Intervalo & Probabilidad \\
  \midrule\noalign{}
  \endhead
  \bottomrule\noalign{}
  \endlastfoot
  \(0 \le x \le 5\) & 0,30 \\
  \(0 \le x \le 7\) & 0,41 \\
  \(0 \le x \le 14\) & 0,30 \\
  \end{longtable}
\end{itemize}

\section{Solution}\label{solution}

Para encontrar la tabla correcta, debemos extraer la información
directamente del gráfico y compararla con las opciones proporcionadas.

\textbf{Paso 1: Analizar la información del gráfico.}

El gráfico está dividido en tres intervalos con sus respectivas
probabilidades:

\begin{enumerate}
\def\labelenumi{\arabic{enumi}.}
\tightlist
\item
  El primer intervalo va desde \(x=0\) hasta \(x=5\). La probabilidad en
  esta sección es \textbf{0,30}.
\item
  El segundo intervalo (central) va desde \(x=5\) hasta \(x=7\). La
  probabilidad en esta sección es \textbf{0,41}.
\item
  El tercer intervalo va desde \(x=7\) hasta \(x=14\). La probabilidad
  en esta sección es \textbf{0,30}.
\end{enumerate}

\textbf{Paso 2: Construir la tabla esperada.}

Con la información anterior, la tabla correcta debe tener la siguiente
estructura:

\begin{longtable}[]{@{}cc@{}}
\toprule\noalign{}
Intervalo & Probabilidad \\
\midrule\noalign{}
\endhead
\bottomrule\noalign{}
\endlastfoot
\(0 \le x \le 5\) & 0,30 \\
\(5 < x \le 7\) & 0,41 \\
\(7 < x \le 14\) & 0,30 \\
\end{longtable}

\textbf{Paso 3: Comparar con las opciones.}

Al revisar las tablas de las opciones, la única que coincide
perfectamente con los datos extraídos del gráfico es la tabla correcta.
Las otras tablas presentan errores comunes:

\begin{itemize}
\tightlist
\item
  \textbf{Error de probabilidad acumulada:} Una de las tablas muestra
  cómo la probabilidad se va sumando (0,30, 0,71, 1,00), lo cual es una
  interpretación incorrecta del gráfico.
\item
  \textbf{Error de definición de intervalo:} Otra tabla usa intervalos
  acumulativos (ej: \(0 \le x \le 7\)), lo cual no corresponde a las
  secciones individuales mostradas.
\item
  \textbf{Error de asignación:} Otra tabla intercambia los valores de
  probabilidad, asignando el valor central a los laterales y viceversa.
\end{itemize}

Por lo tanto, la respuesta correcta es la que representa fielmente los
tres intervalos y sus probabilidades tal como se muestran en el gráfico.

\subsection{Answerlist}\label{answerlist-1}

\begin{itemize}
\item
  \begin{longtable}[]{@{}cc@{}}
  \toprule\noalign{}
  Intervalo & Probabilidad \\
  \midrule\noalign{}
  \endhead
  \bottomrule\noalign{}
  \endlastfoot
  \(0 \le x \le 5\) & 0,41 \\
  \(5 < x \le 7\) & 0,30 \\
  \(7 < x \le 14\) & 0,30 \\
  \end{longtable}
\end{itemize}

\textbf{Incorrecto.} Esta tabla no coincide con la información del
gráfico. Revisa con atención los límites de los intervalos y los valores
de probabilidad asignados a cada sección. * \textbar{} Intervalo
\textbar{} Probabilidad \textbar{}
\textbar:---------------:\textbar:------------:\textbar{} \textbar{}
\(0 \le x \le 5\) \textbar{} 0,30 \textbar{} \textbar{} \(5 < x \le 7\)
\textbar{} 0,71 \textbar{} \textbar{} \(7 < x \le 14\) \textbar{} 1,00
\textbar{}

\textbf{Incorrecto.} Esta tabla no coincide con la información del
gráfico. Revisa con atención los límites de los intervalos y los valores
de probabilidad asignados a cada sección. * \textbar{} Intervalo
\textbar{} Probabilidad \textbar{}
\textbar:---------------:\textbar:------------:\textbar{} \textbar{}
\(0 \le x \le 5\) \textbar{} 0,30 \textbar{} \textbar{} \(5 < x \le 7\)
\textbar{} 0,41 \textbar{} \textbar{} \(7 < x \le 14\) \textbar{} 0,30
\textbar{}

\textbf{Correcto.} Esta tabla representa exactamente los tres intervalos
y sus probabilidades correspondientes como se muestra en el gráfico. *
\textbar{} Intervalo \textbar{} Probabilidad \textbar{}
\textbar:----------------:\textbar:------------:\textbar{} \textbar{}
\(0 \le x \le 5\) \textbar{} 0,30 \textbar{} \textbar{}
\(0 \le x \le 7\) \textbar{} 0,41 \textbar{} \textbar{}
\(0 \le x \le 14\) \textbar{} 0,30 \textbar{}

\textbf{Incorrecto.} Esta tabla no coincide con la información del
gráfico. Revisa con atención los límites de los intervalos y los valores
de probabilidad asignados a cada sección.

\section{Meta-information}\label{meta-information}

exname: probabilidad\_distribucion\_grafico\_tabla extype: schoice
exsolution: 0010 exshuffle: TRUE exsection:
Estadística/Probabilidad/Interpretación de gráficos

\end{document}
